\section{The Night of 8 October 1871}

\subsection{What burned, and where}

On the night of 8 October 1871 --- the same night as the Chicago fire
--- firestorms swept both shores of Green Bay. The compilation
published at Green Bay within weeks by Frank Tilton describes ``the
terrible and unparalleled conflagrations that desolated the North
Eastern counties of Wisconsin on the night of October 8th, 1871, and in
a few short hours swept from Time into Eternity more than a thousand
men, women and children, rendered about four thousand people homeless
and utterly destitute.''\endnote{Frank Tilton, \work{Sketch of the
Great Fires in Wisconsin at Peshtigo, the Sugar Bush, Menekaune,
Williamsonville, and Generally on the Shores of Green Bay, with
Thrilling and Truthful Incidents by Eye Witnesses. Also, a Statement of
Relief Funds Contributed} (Green Bay: Robinson \& Kustermann,
Publishers, 1871), Preface. \textbf{Artifact caution:} the copy read is
not the 1871 imprint. It is a modern reset reproduction, digitized and
uploaded to the Internet Archive in 2022 under the identifier
\texttt{sketch-of-the-great-fires-in-wisconsin}, which reproduces the
1871 title page typographically and includes the folded ``Map of the
Burnt District of Wisconsin 1871'' as a plate. Pages cited were read at
150~dpi. Nothing in this study depends on a reading that a
reset edition could not preserve, and the one place where the reset
text is demonstrably corrupt is flagged below.}

Champion is on the east shore. It was not, strictly, in the Peshtigo
fire. Tilton is careful about this: the east-shore firestorm ``followed
the same general direction of the one on the west shore, while it was
distinct from it, the two being about thirty miles apart, and no
noticeable gale occurring on the bay lying between.'' Pernin says the
same thing from the other side, placing the Robinsonville event ``at
nine or ten leagues distance'' from Peshtigo, ``farther separated from
it by the waters of Green Bay.'' The popular name ``the Peshtigo
fire,'' applied to Champion, names the night and the regional
catastrophe, not the identical flame front.\endnote{Tilton, \work{Sketch
of the Great Fires}, ch.~XIV; Pernin, \work{The Finger of God Is
There!}, p.~94; French, p.~91 (``\`a 9 ou 10 lieues loin de moi; et
cette distance qui nous s\'epare est form\'ee par les eaux de la Baie
Verte''). The reset Tilton text at this point reads ``On the night of
October 3d, at about the same hour of the disaster on the west shore of
Green Bay'' --- a date that contradicts the same volume's own preface,
its map, and every other witness. The reading was checked against the
page image and is what the reset edition prints; whether the error is
the 1871 compositor's or the modern setter's cannot be determined from
this artifact. This study treats the east-shore firestorm as belonging
to the night of 8 October 1871 on the strength of the preface, the
map, and Pernin.}

Tilton's description of the east-shore track is where the Champion
question becomes sharp:

\begin{witnesstext}{Tilton, \work{Sketch of the Great Fires in
Wisconsin} (Green Bay, 1871), ch.~XIV, on the east shore of Green Bay}
The largest settlements destroyed in this track of fire were known as
New Franken, Walhein, Robinsonville, Harris' Pier, Thiry Daems,
Dyckesville, St.~Sauveur, Rosiere, Williamsonville and Forestville.
None of these were villages, but they were either agricultural or
lumbering communities, or settlements around religious institutions.
\dots{} The district burned over here covered an area of fully 500
square miles, and in that extent of country the desolation was almost
complete. If a building escaped it was an exception to the general
rule.\endnote{Tilton, \work{Sketch of the Great Fires}, ch.~XIV, read
at 150~dpi. ``Robinsonville'' also appears on the volume's map of the
burnt district, on the east shore of Green Bay within the shaded area.
The Door County casualty and building figures Tilton gives in the same
paragraph (117 of about 130 deaths; two mills, one church, six
schoolhouses, 148 dwellings burned) are county figures and do not
cover Robinsonville, which is in Brown County.}
\end{witnesstext}

So the earliest account of the east-shore fire that this study located
--- compiled at Green Bay, within weeks, by a writer collecting relief
subscriptions and eyewitness statements --- names Robinsonville among
the settlements destroyed, records no exception there, and does not
mention a chapel, a procession, or a preservation. It says in the same
breath that a surviving building anywhere in the district ``was an
exception to the general rule,'' which shows that Tilton was alert to
exceptions and thought them worth remarking.

That is an argument from silence, and it has the ordinary weaknesses of
one. Tilton was compiling at speed, from the west shore's perspective,
for a relief-fundraising volume; Robinsonville was a Walloon-speaking
Catholic settlement with no newspaper and no obvious channel into his
sources; a chapel and a schoolhouse inside a six-acre fence would not
have altered his casualty tables. His list of ``settlements destroyed''
is a list of places the fire went through, not a building-by-building
survey. The silence is real, and it is not decisive.

\subsection{What Pernin reports}

Three years later, Pernin gives the account that everything since
depends on. It is short, and it is worth having exactly.

\begin{witnesstext}{Pernin, \work{The Finger of God Is There!},
Appendix, printed pp.~99--101 (Montreal, 1874), with the French
collated}
On the 8th of October, 1871, this Belgian colony was visited by the
same whirlwind of wind and fire that overwhelmed Peshtigo, and though
the destruction of human life was less great than in the latter place,
it nevertheless reduced to ashes, farms, houses and timber, covering a
surface several leagues in extent. Now, when the hurricane burst forth,
those pious girls said between each other: ``If the Blessed Virgin
still has need of us, she will protect our lives, if not, we must
succumb to the fiery death awaiting us.'' Filled with confidence and
resignation, they hastened to the chapel, reverently raised the statue
of Our Lady, and kneeling, bore it in procession around their beloved
Sanctuary, reciting their beads. When flame and wind blew so strongly
in the direction of the chapel as to prevent their farther progress,
unless they exposed themselves to suffocation, they awaited a lull in
the storm, or turning in another direction, continued to hope and
pray.\par
Thus passed for them the long hours of that terrible night. I know not
if, supported only by nature, they would have been able to live through
that awful ordeal, but I feel convinced they could not have done it, at
Peshtigo, without a miracle.\par
Morning's light revealed the deplorable ravages wrought by the
conflagration. All the houses and fences in the neighborhood had been
burned, with the exception of the school, the chapel and fence
surrounding the six acres of land consecrated to the Blessed Virgin.
This paling had been charred in several places, but the fire, as if it
had been a sentient being, whilst consuming everything in the vicinity,
the winding path surrounding the enclosure being only eight or ten feet
wide, had respected this spot, sanctified by the visible presence of
the Mother of God, and it now shone out like an emerald island amid a
sea of ashes.\endnote{Pernin, \work{The Finger of God Is There!},
pp.~99--101; French, pp.~97--99. The English scan's optical-character
recognition is badly damaged through the sentence beginning ``If the
Blessed Virgin still has need of us'' and again at the close of the
passage; both were therefore read directly at 200~dpi from the page
images of printed pp.~99--100, which carry the wording given here. The
French reads at the first place ``\emph{Si la Ste.\ Vierge a besoin
encore de nous, elle saura nous prot\'eger, si non, elle nous laissera
br\^uler avec les autres},'' a slightly blunter formulation than the
English. The
French for the last sentence reads: ``\emph{Cette place comme
sanctifi\'ee par l'attouchement de la M\`ere de Dieu, apparaissait
semblable \`a une \^ile verdoyante au milieu d'une mer de
cendres.}''}
\end{witnesstext}

Pernin then adds, under the heading ``Important consideration'':

\begin{witnesstext}{Pernin, \work{The Finger of God Is There!}, printed
p.~101, ``Important consideration''}
In relating the preceding fact I have no intention of pronouncing it a
miracle, no more than I would adventure to qualify as miracle the
preservation of the tabernacle in the midst of the fires of Peshtigo.
These two facts greatly edified myself, enlivening my faith and hope,
and in narrating them I have no other aim than that of edifying others.
I have no intention either of passing judgment on the apparition of the
Blessed Virgin and on the pious pilgrimages which have resulted from
it. Ecclesiastical authority has not as yet spoken on the subject; it
silently allows the good work to advance, awaiting perhaps some proof
more striking and irrefutable before pronouncing its fiat. Far from me
be the thought of forestalling ecclesiastical judgment.\endnote{Pernin,
\work{The Finger of God Is There!}, p.~101; French, pp.~99--100
(``\emph{R\'eflexion importante}''). Pernin's parallel refusal about
his own tabernacle at Peshtigo is in the same sentence in both
languages.}
\end{witnesstext}

Everything in that account, and nothing beyond it, is what Pernin
witnesses. Note in particular who is present: ``those pious girls,''
\term{ces pieuses filles} --- the five or six young women who had
joined Adele, introduced two pages earlier. Pernin describes a small
community praying round its own enclosure. He does not describe a
refugee crowd, livestock, a well, or rain. He does not say the fire was
extinguished; he says that morning came and the enclosure was intact.
And he twice declines, in terms, to call any of it a miracle or to
anticipate the judgment of the Church.

\subsection{A quotation that is not Pernin's}

The shrine's current narrative of the fire introduces a paragraph with
the words: ``The following is the account of Father Peter Pernin, a
local priest who described the grounds after the events.'' The
paragraph runs:

\begin{witnesstext}{Paragraph attributed to ``Fr.\ Peter Pernin'' on
the National Shrine's ``Our Story'' page, read 2026-07-25}
After hours of horror and suspense, the heavens sent relief in the form
of a downpour. The fervent prayers to the Mother of God were heard. The
fire was extinguished, but dawn revealed the ravages wrought by the
conflagration. Everything about them was destroyed; miles of desolation
everywhere. But the convent, school, and chapel on the holy land
consecrated to the Virgin Mary shone like an emerald isle in a sea of
ashes. The raging fire licked the outside palings and left charred
scars as mementos. Tongues of fire had reached the chapel fence, and
threatened destruction to all within its confines; the fire had not
entered the Chapel grounds.\endnote{``Our Story,''
\url{https://championshrine.org/our-story/}, read 2026-07-25, under the
heading ``The Fire,'' attributed on the page to ``Fr.~Peter Pernin.''}
\end{witnesstext}

This study collated that paragraph against both 1874 printings of
Pernin's book, English and French, in full. It is not his text. The
resemblance is confined to one image --- the emerald isle in a sea of
ashes --- and to the charred palings. Everything that carries weight in
the paragraph is absent from Pernin:

\begin{boundarytable}{In the paragraph attributed to Pernin}{In Pernin,
1874, English and French}
``After hours of horror and suspense, the heavens sent relief in the
form of a downpour.'' & No rain of any kind, in either language.
Pernin's transition is simply ``Morning's light revealed'' /
``\emph{Le matin venu, on constata}.'' \\
``The fire was extinguished.'' & Pernin never says the fire was
extinguished; he reports what daylight showed. \\
``The convent, school, and chapel.'' & ``The school, the chapel and
fence surrounding the six acres.'' No convent. \\
``The fervent prayers to the Mother of God were heard.'' & Pernin
declines to draw the conclusion: ``I have no intention of pronouncing
it a miracle.'' \\
``Everything about them was destroyed; miles of desolation
everywhere.'' & ``All the houses and fences in the neighborhood had
been burned,'' with the enclosure excepted. Comparable in substance. \\
``Shone like an emerald isle in a sea of ashes''; the charred palings &
``Shone out like an emerald island amid a sea of ashes''; ``this paling
had been charred in several places.'' The only genuine verbal
correspondence. \\
\end{boundarytable}

This study does not know where the paragraph comes from, and says so.
It reads as a later retelling of Pernin --- possibly the shrine's own
mid-twentieth-century history, which this study could not reach ---
that has acquired his name in transmission. What can be said with
confidence is narrower and sufficient: the rain, the extinguishing of
the fire, the answered prayer, and the convent are not in the 1874
text, and any argument that rests on Pernin's authority for them rests
on nothing.

\subsection{What the 2010 decree says about the fire}

The decree includes the fire among its recitals, in one sentence:

\begin{witnesstext}{Decree of 8 December 2010, p.~2, recital}
This holy place was preserved from the infamous Peshtigo fire of 1871,
when many of the faithful gathered here with Adele and prayed through
the intercession of Our Lady of Good Help, with the result that the
fire that devastated everything in its wake in this entire area stopped
when it reached the parameters of the Shrine.\endnote{Decree of 8
December 2010, p.~2. The recital adopts the fuller narration --- ``many
of the faithful'' rather than Adele's small community. It stands among
the \textsc{Given that} recitals, not in the \textsc{Therefore} clause.}
\end{witnesstext}

Two things follow, and they are easy to conflate. The fire is part of
the decree's \emph{reasoning}: it belongs to the recitals that
establish a continuous history of faith and devotion at the site. It is
not part of the decree's \emph{object}: the operative clause judges the
events, apparitions, and locutions of October 1859, and nothing else.
The decree's only explicit statement about miracles cuts the other way
and covers the reported cures: ``though none of these favors have been
officially declared a miracle by the Church, they are clear evidence of
spiritual fruitfulness.''\endnote{Decree of 8 December 2010, pp.~1--2.}
No competent authority has declared the preservation of 8 October 1871
a miracle, and the decree does not do so.

\begin{judgmentbox}{Project synthesis --- what the fire evidence will
carry}
This study judges that the record supports the following and no more.
\textbf{Carried:} that the east-shore firestorm of the night of 8
October 1871 passed through the Robinsonville settlement and destroyed
it, on the testimony of a Green Bay compilation of the same year;
that a small community of women spent that night praying and
processing with a statue around a fenced six-acre enclosure; and that
the next morning the school, the chapel, and the fence stood, charred
but not consumed, amid general destruction, on the testimony of a
priest writing three years later who had access to the participants and
who published his account under an episcopal recommendation.
\textbf{Not carried:} rain; the extinguishing of the fire; a refugee
crowd; livestock at the chapel; the shallow well; and any of these on
Pernin's authority. \textbf{Not adjudicated by anyone:} whether the
preservation was miraculous. Pernin declined to say; the 2010 decree
recites the fact and does not declare it; no other competent act was
located.\jsep
The strongest counterargument is that Tilton's silence proves nothing
and that oral tradition in a small Walloon settlement could preserve
the rain, the crowd, and the cattle perfectly well for eighty years
before anybody wrote them down --- and that survivors' descendants
still keep the anniversary, which is itself a form of transmission.
That is fair, and it is why this study classifies those elements as
undocumented rather than as false. What would move them into the
carried column is any dated record before 1955 --- a diocesan or parish
note, a newspaper account, a letter, a claim in the fire-relief papers
--- that reports the crowd, the animals, or the rain. What would move
them out is a demonstration that they first appear in a
mid-twentieth-century retelling.
\end{judgmentbox}

\subsection{A pastoral note that belongs beside the claim}

The east-shore firestorm killed people. Tilton counts about 130 dead in
the peninsula track and 117 of them in Door County alone; the two
shores together took more than a thousand lives. A preservation
narrative told without that fact becomes something the sources
themselves refuse: an implied verdict on the faith of the dead. Pernin,
who was in the river at Peshtigo while his own parishioners burned,
never draws it. Neither does the 2010 decree. Neither may anyone else.
Whatever happened inside that fence, it is not evidence about the souls
outside it.
